\section{Fase 3: Integracion LLM y n8n}
\label{sec:integracion-llm-n8n}

Con la fase 2 terminada y viendo que las limitaciones eran evidentes, decidimos hacer un cambio muy importante para el futuro del proyecto. Abandonamos 
la idea del sistema de reglas que teniamos, aunque no en la totalidad, ya que había partes que nos seguirían siendo útiles. Decidimos así integrar un
modelo de lenguaje (LLM) como nucleo o motor del análisis y generación de las paginas. Hasta este momento el sistema intentaba entender los datos mediante
puntuaciones y comparaciones, pero decidimos que fuese la Inteligencia Artificial la encargada de escoger y proponer la estructura del sitio y decidiera que 
campos mostrar. Como comente antes, no todo el flujo de análisis e ingesta se cambio, se eliminaron las partes de score y propuesta de dataset, pero en parseo 
y análisis de los campos se mantuvo y modifico levemente, se metió todo el dataset en una estructura, que seria la que se pasaria al LLM, de forma que evitariamos
alucinaciones o errores del LLM en la lectura de la API. Esta parte esta mejor explicada en la subseccion~\ref{subsec:proteccion-anti-alucionaciones}


\subsection{Conexion con LLM}
\label{subsec:conexion-llm}

Para poder realizar la conexion con el LLM, el primer paso fue crear el módulo \textbf{utils/llm}, y dentro creamos el archivo \texttt{client.py}, el cual
implementa la funcion \texttt{chat\_completions}. Esta funcion se encarga de hacer una unica cosa: enviar un mensaje al LLM y devolver la respuesta en texto
plano.

Una decisión importante fue utilizar el estandar de la mayoria de proveedores para realizar la conexion, utilizando de esta forma el endpoint \textbf{/chat/completions}. En este caso creado para poder realizar pruebas con varios LLM de diferentes proveedores (OpenRouter, Groq, NVIDIA) y cualquier agente local.
Otra decision fue crear \texttt{.env}, de forma que la URL base, la API key y el modelo se leen de variables de entorno definidas en dicho fichero, lo que permite cambiar de proveedor sin tocar una sola línea de código.

La función mencionada anteriormente acepta tres parámetros principales:

\begin{itemize}
    \item \textbf{Mensaje de usuario}: la tarea concreta que se le envía al modelo 
    en cada llamada, en este caso el prompt con los datos del dataset, las claves 
    disponibles y los ejemplos de items reales.
    
    \item \textbf{Mensaje de sistema}: instrucciones previas que definen cómo debe 
    comportarse el modelo antes de procesar el mensaje del usuario. Se utiliza para 
    indicarle el rol que debe adoptar, el formato de respuesta esperado y lo que 
    está prohibido hacer. Es opcional, ya que técnicamente todo puede ir en el 
    mensaje de usuario, pero separarlo es una buena práctica que facilita reutilizar 
    el mismo comportamiento base en distintas llamadas.
    
    \item \textbf{Temperatura}: parámetro que controla el grado de aleatoriedad 
    en las respuestas del modelo. Un valor cercano a 0 produce respuestas más 
    deterministas y predecibles, mientras que valores más altos generan respuestas 
    más creativas pero menos consistentes. En este caso se fija por defecto en 
    \texttt{0.2}, favoreciendo la estabilidad especialmente cuando se espera JSON 
    estructurado como salida.
\end{itemize}

El primer proveedor utilizado fue OpenRouter, con el modelo de meta-llama: \\
\texttt{llama-3.3-70b-instruct:free}, de acceso gratuito y potencia más que suficiente para las primeras pruebas.


\subsection{Planner schema}
\label{subsec:planner-schema}

Otro fichero clave fue \texttt{utils/llm/planner.py} y su funcion \texttt{generate\_site\_plan}. La mision era transformar el resultado del analisis en un propuesta con campos
concretos de schema que el usuario pudiese aceptar y modificar a su gusto antes de continuar.

El schema tiene una forma fija que el resto del sistema espera:
\begin{itemize}
    \item \textbf{site\_type}: tipo de sitio recomendado, entre \texttt{blog}, \texttt{portfolio}, \texttt{catalog}, \texttt{dashboard} u \texttt{other}.
    \item \textbf{site\_title}: nombre descriptivo del sitio propuesto por el modelo.
    \item \textbf{fields}: lista ordenada de campos del dataset que el sitio debe mostrar, cada uno con su clave original y una etiqueta legible para el usuario.
\end{itemize}

Para obtener este esquema, se construye un prompt estructurado que se envía al LLM. Este prompt es el resultado de una serie de pruebas con distintas formulaciones,  
siguiendo técnicas de \emph{prompt engineering}~\cite{openai_prompt_engineering} para conseguir respuestas consistentes y parseables. La información contenida en este 
prompt es la siguiente:

\begin{itemize}
    \item \textbf{Objetivo del usuario}: texto en lenguaje natural introducido por 
    el usuario describiendo qué tipo de sitio quiere construir.
    
    \item \textbf{Claves disponibles}: lista de campos reales extraídos del dataset 
    durante la fase de análisis. El modelo solo puede proponer campos que estén 
    en esta lista.
    
    \item \textbf{Ejemplos del dataset}: hasta tres items reales del dataset, para 
    que el modelo pueda entender qué tipo de valores contiene cada campo y proponer 
    etiquetas adecuadas.
    
    \item \textbf{Contrato JSON}: estructura exacta que debe devolver el modelo, 
    con los campos obligatorios y su formato esperado.
    
    \item \textbf{Reglas de formato}: conjunto de instrucciones explícitas que 
    indican al modelo que debe devolver únicamente un objeto JSON válido, sin 
    bloques de código Markdown, sin texto introductorio y sin ningún contenido 
    fuera del JSON.
\end{itemize}

Esta rigidez sobre el formato de salida, es intencionada, ya que tras la realizacion de varias pruebas los fallos generados por los diferentes LLM siempre eran parecidos.
Como esa respuesta va a ser usada directamente, cualquier texto adicional rompe el proceso. De todas formas, debido a las alucinaciones de algunos modelos y la ignoracion 
de las restricciones en algunas ocasiones, se aplicaba una limpieza previa que eliminaba bloques Markdown y caracteres sobrantes del parseo JSON.


\subsection{Proteccion anti-alucinaciones}
\label{subsec:proteccion-anti-alucionaciones}

Los errores mas comunes siempre sucedian por la misma causa, las alucinaciones del LLM. El error mas comun era inventar informacion que no existia en la API, y que nunca se
en el parseo, por lo que no existe en el dataset contruido. Para ello hubo que aplicar varias capas de defensa contra alucinaciones para proteger los schemas.

\begin{itemize}
    \item \textbf{Validación contra claves reales}: cada clave propuesta por el LLM se comprueba contra la lista de claves reales extraída del dataset. Si la clave no existe, se intenta un match sin distinguir mayúsculas y minúsculas. Si tampoco hay coincidencia, la clave se descarta silenciosamente.
    \item \textbf{Filtro de tokens prohibidos}: existe una lista de palabras que el modelo confunde con claves del dataset porque forman parte del propio prompt, como \texttt{available\_keys}, \texttt{fields} o \texttt{examples}. Cualquier clave que coincida con esta lista se descarta automáticamente.
    \item \textbf{Detección de valores disfrazados de claves}: se detectan los casos en los que el modelo introduce un valor concreto del dataset donde debería haber puesto el nombre de un campo, como fechas en formato ISO, URLs completas o cadenas excesivamente largas.
    \item \textbf{Fallback}: si tras aplicar todas las validaciones no queda ninguna clave válida, el sistema construye automáticamente un schema de emergencia con las primeras claves del dataset, garantizando que el flujo nunca se bloquea por un fallo del modelo.
\end{itemize}

Si el parseo falla, el sistema reintenta automáticamente una vez con una pista sobre el error cometido. Si el segundo intento también falla, se activa el fallback.


\subsection{Actualizacion de modelos}
\label{subsec:actualizacion-modelos}

La integracion del LLM trajo consigo dos cambios importantes.

El primero de ellos fue una modificacion en el modelo existente \textbf{APIRequest}, añadiendole el campo \texttt{plan\_accepted}. Siendo este un booleano encargado de guardar
si el usuario ha aceptado el schema propuesto antes de continuar.

El segundo fue la creacion de un nuevo modelo llamado \textbf{GeneratedSite}, vinculado con el modelo APIRequest mediante una relacion \texttt{OneToOne}. Este modelo se
encargaba de guardar el plan aceptado, y los templates generados para el sitio web, y un identificador unico en el campo \texttt{public\_id}. Este modelo fue muy importante para
la generacion de los proyectos que se desarrollaria en la siguiente fase.


\subsection{Integracion de n8n}
\label{subsec:integracion-n8n}

En paralelo a la integración del LLM, durante esta fase se incorporó n8n al sistema por primera vez. Es una herramienta de automatización de flujos que permite conectar servicios mediante webhooks y nodos, y que más adelante sería una pieza fundamental de la arquitectura del proyecto. En este momento su uso fue todavía inicial: instalar el entorno y crear el primer flujo de pruebas.

n8n se instaló en un contenedor Docker siguiendo la imagen oficial, quedando disponible en \texttt{localhost:5678}. Una vez levantado el entorno, se creó el primer flujo en n8n y se configuró el primer webhook, apuntando a la ruta \texttt{/webhook-test/WebBuilder-Register}. El prefijo \texttt{webhook-test} indica que el webhook estaba en modo de prueba, el modo que n8n ofrece para verificar que el flujo funciona correctamente antes de activarlo en producción.

La integración en Django se realizó en \texttt{views/auth.py}, donde se añadió la función \texttt{llamar\_webhook}.
El nodo se encargaría de notificar a un nuevo usuario en base de datos, es decir, aquellos usuarios recién registrados. Django realizaba una llamada POST a ese webhook con el nombre de usuario y el email, y n8n enviaba automáticamente un correo de bienvenida. Más adelante se completaría el flujo con el webhook de login.

La función estaba diseñada para ser completamente transparente al flujo principal: si n8n no estaba disponible o la llamada fallaba, la excepción se capturaba silenciosamente y el registro continuaba con normalidad. Esto garantizaba que un fallo en el sistema de notificaciones nunca pudiera bloquear el acceso a la aplicación.


